\documentclass[tikz,border=5pt]{standalone}
\usepackage[UTF8,fontset=fandol]{ctex}
\usepackage{amsmath,amssymb}
\usetikzlibrary{arrows.meta,calc,positioning}
\definecolor{ink}{HTML}{18324A}
\definecolor{blue}{HTML}{247DA6}
\definecolor{orange}{HTML}{DD8A28}
\definecolor{green}{HTML}{27896C}
\begin{document}
\begin{tikzpicture}[>=Latex,font=\small,text=ink,arr/.style={->,line width=1pt,blue},box/.style={draw=blue,fill=blue!5,rounded corners=2pt,align=center,font=\footnotesize,minimum height=1.1cm,text width=3.0cm}]
\path[use as bounding box] (-.1,-5.6) rectangle (16.1,3.2);
\node[anchor=west,font=\bfseries] at (0,2.9) {(a) 6D = 3维平移 + 3维旋转};
\begin{scope}[shift={(1,.5)}]
\draw[fill=blue!8,draw=blue,thick] (0,0)--(1.8,0)--(1.8,1.1)--(0,1.1)--cycle;
\draw[fill=blue!15,draw=blue,thick] (0,1.1)--(.6,1.55)--(2.4,1.55)--(1.8,1.1)--cycle;
\draw[fill=blue!25,draw=blue,thick] (1.8,0)--(2.4,.45)--(2.4,1.55)--(1.8,1.1)--cycle;
\draw[->,red!70!black,thick] (0,0)--(1.4,0) node[below] {$x_o$};
\draw[->,green,thick] (0,0)--(0,1.5) node[above] {$y_o$};
\draw[->,blue,thick] (0,0)--(.6,.45) node[right] {$z_o$};
\node[align=center] at (1,-.6) {物体局部坐标\\模型／可见表面参考};
\end{scope}
\draw[arr,line width=1.5pt] (4,.9)--node[above,align=center] {$R\in SO(3)$，$t\in\mathbb R^3$\\旋转并平移}(7,.9);
\node[align=center] at (5.5,-.35) {$x_c=Rx_o+t$\\$T_{oc}=\begin{bmatrix}R&t\\0&1\end{bmatrix}$};
\begin{scope}[shift={(7.8,.25)}]
\draw[->,red!70!black,thick] (0,0)--(2,0) node[right] {$x_c$};
\draw[->,green,thick] (0,0)--(0,1.6) node[above] {$y_c$};
\draw[->,blue,thick] (0,0)--(-.55,.6) node[above] {$z_c$};
\draw[arr,draw=orange] (0,0)--(1.1,.65) node[midway,above left] {$t$};
\begin{scope}[shift={(1.1,.65)},rotate=18,scale=.6]
\draw[fill=orange!12,draw=orange,thick] (0,0)--(1.8,0)--(1.8,1.1)--(0,1.1)--cycle;
\draw[fill=orange!20,draw=orange,thick] (0,1.1)--(.6,1.55)--(2.4,1.55)--(1.8,1.1)--cycle;
\draw[fill=orange!30,draw=orange,thick] (1.8,0)--(2.4,.45)--(2.4,1.55)--(1.8,1.1)--cycle;
\end{scope}
\node[align=center] at (1.05,-.65) {相机坐标中的物体\\形状保持刚性};
\end{scope}
\node[box,text width=3.3cm] (pixel) at (14,.75) {像素投影\\$u=\pi(Kx_c)$\\框、坐标轴、模型轮廓};
\draw[arr] (11.05,.75)--(pixel);
\node[anchor=west,font=\bfseries] at (0,-1.5) {(b) FoundationPose：首帧注册与后续跟踪};
\node[box] (input) at (1.6,-2.55) {RGB + 米制深度\\内参 $K$ + 目标mask\\对应mesh};
\node[box] (hyp) at (5.6,-2.55) {初始化姿态候选\\学习式细化\\学习式评分};
\node[box] (pose) at (9.6,-2.55) {首帧注册位姿\\$T_{oc}^{(0)}$\\绑定目标身份};
\node[box] (track) at (13.6,-2.55) {前帧位姿 + RGB-D\\当前帧连续跟踪\\$T_{oc}^{(t)}$};
\draw[arr] (input)--(hyp);\draw[arr] (hyp)--(pose);\draw[arr] (pose)--(track);
\draw[arr] (track.east)--(15.5,-2.55)--(15.5,-4)--(12.6,-4)--(12.6,-3.12);
\node[font=\scriptsize,fill=white] at (14.05,-4) {传递到下一帧};
\node[draw=green,fill=green!5,rounded corners,align=center,text width=9.8cm] at (5.1,-4.25) {输出检查：同帧／同物体／同单位；$R^{\mathsf T}R\approx I$，$\det R\approx1$\\独立姿态精度需同坐标真值与对称性处理，不能只看框贴合};
\node[anchor=west,font=\scriptsize] at (0,-5.2) {原理示意：方块仅代表刚体，不是实验mesh；旋转六维编码是特征表示，不等于这里的6个自由度。};
\end{tikzpicture}
\end{document}
